% U7 WS2 -- magnitude and properties of K. Block 3.
\documentclass{apchem-worksheet}

\apcunit{7}
\apcunittitle{Equilibrium}
\apcdoctitle{Worksheet 2 \textbullet\ Magnitude and Properties of $K$}
\apcsubtitle{CED 7.5--7.6 \textbullet\ equilibrium constants multiply where enthalpies add}

\begin{document}
\wsheader[Reverse $\Rightarrow$ $1/K$ \textbullet\ multiply coefficients by
$n$ $\Rightarrow$ $K^n$ \textbullet\ divide by $n$ $\Rightarrow$
$K^{1/n}$ \textbullet\ add two reactions $\Rightarrow$ $K_1K_2$
\textbullet\ only temperature changes $K$.]

\problem[]{\ced{7.5} Classify each as products-favored,
reactants-favored, or appreciable amounts of both:}
\begin{parts}
  \item $K = 1\times10^{-12}$ \answerblank[5.4cm]{reactants strongly favored}
  \item $K = 2.5$ \answerblank[5.6cm]{appreciable amounts of both}
  \item $K = 4\times10^{8}$ \answerblank[5.2cm]{products strongly favored}
  \item $K = 3\times10^{-3}$ \answerblank[5cm]{reactants favored}
\end{parts}

\problem[]{\ced{7.5} A reaction has $K = 10^{35}$, yet mixing the reactants
produces no observable change over several hours.}
\begin{parts}
  \item Is the value of $K$ necessarily wrong? \answerblank[1.8cm]{no}
  \item Explain. \answerlines[3]{$K$ describes only the position of
        equilibrium, not the rate of getting there. A very high activation
        energy can make the reaction immeasurably slow, so the system is
        under kinetic control while the thermodynamic prediction remains
        correct.}
  \item What single change would speed it up without altering $K$?
        \answerblank[3.6cm]{adding a catalyst}
\end{parts}

\problem[]{\ced{7.6} For a reaction with $K = 4.0$, give $K$ for:}
\begin{parts}
  \item the reverse reaction \answerblank[2.4cm]{0.25}
  \item the reaction with all coefficients doubled
        \answerblank[2.4cm]{16}
  \item the reaction with all coefficients halved
        \answerblank[2.4cm]{2.0}
  \item this reaction added to a second one with $K = 3.0$
        \answerblank[2.4cm]{12}
\end{parts}

\problem[]{\ced{7.6} In Unit 6 you learned that combining reactions
\emph{adds} their $\Delta H$ values. Here, combining reactions
\emph{multiplies} their $K$ values. Explain why the two behave differently.
\answerlines[4]{Enthalpy is an additive quantity --- energies of successive
steps sum. An equilibrium constant is a ratio of concentration terms raised
to powers, and combining reactions multiplies those ratios together. The
underlying link is that $\Delta G^\circ$, which adds, is proportional to
$\ln K$ --- and adding logarithms corresponds to multiplying the constants
themselves.}}

\problem[]{\ced{7.6} $K = 2.5\times10^{-3}$ for
$\ce{A + B <=> C}$ at a given temperature.}
\begin{parts}
  \item $K$ for $\ce{C <=> A + B}$: \answerblank[3cm]{$4.0\times10^{2}$}
  \item $K$ for $\ce{2A + 2B <=> 2C}$: \answerblank[3cm]{$6.3\times10^{-6}$}
  \item Which side is favored for the original reaction?
        \answerblank[2.6cm]{reactants}
  \item Which side is favored for the reverse?
        \answerblank[2.6cm]{products}
\end{parts}

\problem[]{\ced{7.6} A student claims that doubling the concentration of a
reactant doubles $K$. Explain what actually happens.
\answerlines[4]{$K$ does not change. Doubling a reactant concentration
changes $Q$, making it smaller than $K$, so the system shifts toward
products until $Q$ returns to the same $K$. The concentrations at the new
equilibrium differ from the old ones, but their combination in the $K$
expression gives the identical value. Only a change in temperature alters
$K$ itself.}}

\problem[]{\ced{7.5} Two reactions at the same temperature have
$K_1 = 5\times10^{-4}$ and $K_2 = 5\times10^{4}$.}
\begin{parts}
  \item Which lies further toward products?
        \answerblank[2.2cm]{the second}
  \item If the second reaction is reversed, how does its $K$ compare with
        $K_1$? \answerlines[2]{Reversing gives $1/(5\times10^{4}) =
        2\times10^{-5}$, which is smaller than $K_1$ --- so the reversed
        second reaction lies even further toward its reactants than the
        first reaction does.}
  \item Does either value tell you which reaction is faster?
        \answerblank[1.8cm]{no}
\end{parts}

\begin{spiralstrip}{Unit 7 \textbullet\ blocks 1--2}
  \item Equilibrium means equal: \answerblank[2cm]{rates}
  \item Pure solids appear in $K$? \answerblank[1.6cm]{no}
  \item $Q < K$ means the system shifts:
        \answerblank[1.8cm]{right}
  \item $K_p = K_c(RT)^{?}$: \answerblank[1.8cm]{$\Delta n$}
  \item $K$ for $\ce{CaCO3(s) <=> CaO(s) + CO2(g)}$:
        \answerblank[2.4cm]{$[\ce{CO2}]$}
\end{spiralstrip}

\end{document}
